% =====================================================================
%  Conclusion générale
% =====================================================================
\chapter*{Conclusion générale}
\addcontentsline{toc}{chapter}{Conclusion générale}
\markboth{Conclusion générale}{}
\label{chap:conclusion}

\section*{Bilan du travail réalisé}

Le projet a abouti à une plateforme fonctionnelle couvrant le parcours complet de
l'adoption de l'intelligence artificielle dans une organisation pharmaceutique, du
diagnostic de maturité à la mesure de la valeur captée.

En chiffres, la réalisation représente 34~tables, 99~endpoints répartis en dix domaines
fonctionnels, sept services métier isolés du transport HTTP, trois migrations de schéma,
quatorze écrans, 41~permissions, onze rôles, cinq niveaux de périmètre et 112~tests
automatisés qui passent.

Neuf des douze modules du cahier des charges sont pleinement implémentés. Trois le sont
partiellement, et le chapitre~\ref{chap:analyse} en détaille l'écart sans le minimiser :
l'export bureautique de la fiche projet produit du HTML imprimable plutôt qu'un fichier
\code{.docx}, le suivi de projet dispose de son API mais pas de son écran dédié, et les
notifications sont persistées sans panneau de consultation.

Trois choix de conception structurent l'ensemble.

\textbf{L'explicabilité comme contrainte de départ, non comme option.} Chaque calcul
renvoie sa décomposition : le readiness rend ses dimensions avec leur preuve, le score
de priorité rend la contribution de ses six critères, la recommandation rend ses règles
satisfaites, ses alternatives écartées et ses garde-fous. Cette contrainte a orienté le
modèle de données — les explications sont persistées en JSON — autant que les formules.

\textbf{La sécurité côté serveur, exclusivement.} L'interface adapte ce qu'elle affiche
aux droits de l'utilisateur, mais aucune décision d'accès ne lui est déléguée. Les
tests d'autorisation appellent l'API directement, sans passer par l'interface, parce que
c'est ainsi qu'un attaquant procéderait.

\textbf{La séparation entre calcul et exposition.} Aucune formule ne vit dans un router.
Cette discipline a rendu les algorithmes testables sans serveur ni base de données, et
c'est ce qui a permis de les corriger tard dans le projet sans crainte de régression.

\section*{Limites de la version actuelle}

Les limites ci-dessous sont réelles et constatées dans le code livré.

\textbf{Aucun modèle d'apprentissage.} La plateforme recommande des solutions
d'intelligence artificielle sans en contenir aucune. Le moteur de qualification est un
système de règles pondérées : il ne découvre rien qu'un expert n'ait prévu et n'apprend
pas des décisions passées. Ce choix est argumenté à la section~\ref{sec:pas-dia} ; il
n'en reste pas moins une limite.

\textbf{La détection de doublons est lexicale.} Le croisement interservices repère les
termes partagés entre deux expressions de besoin. Deux besoins identiques formulés avec
des vocabulaires métier différents ne sont pas rapprochés.

\textbf{Le suivi de projet n'a pas d'écran.} Phases, jalons, tâches, risques,
mitigations, incidents et passage de gates sont pilotables par l'API, pas par
l'interface. Le PMO retrouve l'avancement et la prochaine gate dans son tableau de bord,
ce qui couvre le pilotage mais pas la gestion détaillée.

\textbf{L'export bureautique n'existe pas.} Le libellé du bouton est en avance sur la
fonction.

\textbf{Pas de réinitialisation de mot de passe autonome.} Un utilisateur qui oublie son
mot de passe doit s'adresser à l'administrateur de son organisation.

\textbf{Pas de rafraîchissement de jeton.} Le jeton est valable douze heures, après quoi
l'utilisateur se reconnecte. Il n'existe ni jeton de rafraîchissement, ni liste de
révocation : un jeton volé reste exploitable jusqu'à son expiration. La seule parade
actuelle est la désactivation du compte, qui est bien effective à chaque requête — le
serveur vérifie l'état du compte après avoir validé le jeton — mais qui suppose de
désactiver l'utilisateur légitime.

\textbf{Configuration par défaut à durcir.} La clé de signature et le mot de passe du
superutilisateur initial portent des valeurs par défaut qui doivent impérativement être
redéfinies avant toute mise en production.

\textbf{Écart d'environnement.} Le développement s'est fait sur PostgreSQL~12.4 en
local, tandis que le fichier Docker déclare PostgreSQL~16. Aucune fonctionnalité
spécifique à la version~16 n'est utilisée, mais l'écart n'a pas été validé en
conditions réelles.

\section*{Perspectives}

Les évolutions ci-dessous sont des propositions. Aucune n'est implémentée.

\textbf{Assistance par modèle de langage, en entrée du moteur.} Un modèle pourrait
pré-remplir les quinze faits de qualification à partir d'une description libre du
besoin, rédiger un premier jet de fiche projet et détecter des doublons sémantiques
là où la comparaison lexicale échoue. Le principe directeur serait de laisser le modèle
\emph{proposer des entrées} et le moteur de règles \emph{décider} : la traçabilité de la
décision resterait intacte, et le gain se situerait sur la saisie, qui est aujourd'hui
le point de friction.

\textbf{Apprentissage sur l'historique des décisions.} Après plusieurs dizaines de
cycles complets — qualification, décision de comité, résultat mesuré — l'organisation
disposerait enfin d'un corpus. Les pondérations de priorisation pourraient alors être
calibrées sur les projets qui ont réellement tenu leurs promesses, plutôt que fixées
a~priori.

\textbf{Écran de pilotage de projet.} Exposer l'API de suivi déjà en place :
diagramme de Gantt des jalons, tableau des tâches, matrice des risques et parcours de
gates.

\textbf{Export bureautique réel.} Générer un fichier \code{.docx} ou un PDF mis en page
côté serveur, pour aligner la fonction sur le libellé.

\textbf{Notifications actionnables.} Panneau de consultation, marquage lu et non lu, et
déclenchement sur les événements qui appellent une action : gate à franchir, indicateur
qui décroche, décision de comité en attente.

\textbf{Durcissement de l'authentification.} Jeton de rafraîchissement, révocation côté
serveur, réinitialisation autonome du mot de passe et, en périmètre réglementé,
authentification à deux facteurs.

\textbf{Exploitation des périmètres par site.} Le modèle porte déjà le rattachement des
utilisateurs aux sites et le niveau de périmètre correspondant. Les filtres appliqués
aujourd'hui s'arrêtent au département ; étendre le bornage au site serait une extension
naturelle du dispositif existant.

\section*{Apport personnel}

Ce projet a demandé de tenir ensemble trois exigences qui se contredisent
régulièrement : la richesse fonctionnelle, la rigueur du contrôle d'accès et
l'honnêteté sur ce que le système fait réellement.

La dernière est celle qui a le plus appris. Il aurait été facile de présenter le moteur
de qualification comme une intelligence artificielle : le vocabulaire du domaine s'y
prête, et personne n'aurait vérifié. Choisir de le nommer pour ce qu'il est — un système
expert à base de règles — et d'argumenter pourquoi ce choix est adapté au problème posé
s'est révélé plus solide que l'inverse. Un moteur explicable se défend devant un comité,
devant un inspecteur et devant un jury ; un modèle qu'on ne sait pas justifier ne se
défend nulle part.
